%! Author = sriadityavedantam
%! Date = 3/31/26

\section{Rational Curves}

Let $C$ be an affine plane curve in $\aa^2$.

\begin{definition}[Rational]
    $C$ is rational if it can be parametrized by rational functions such that there exists rational functions $\varphi,\psi$ such that
    \begin{align*}
        x &= \varphi(t)\\
        y &= \psi(t),
    \end{align*}
    and $f(x,y) = 0$.
\end{definition}

\begin{definition}[Rational Function]
    A quotient of rational polynomials.
\end{definition}

\[
\frac{t^3 + 3t^2 + 5}{t - 1}
\]
is a rational function.
The domain is $\c\setminus\{1\}$.
In fact, any complex line is rational due to parametrization.

\[
f(x,y) = 2x + 3y - 5,
\]
We can parameterize
\begin{align*}
    x &= t\\
    y &= \frac{-2x + 5}{3}
\end{align*}

\[
f(x,y) = x - 2,
\]
\begin{align*}
    x &= 2\\
    y &= t
\end{align*}

We can even parametrize the conic $C: \{y - x^2 = 0\}$ as
\begin{align*}
    x &= t\\
    y &= t^2
\end{align*}

But simply allowing one to keep parametrizing these conics through a process of implicit isolation, we cannot have a rational function at all times.
For example:

\[
C: \{x^2 - y^2 = 1\}
\]
\begin{align*}
    x &= t\\
    y &= \sqrt{t^2 - 1},
\end{align*}
but the square root is not rational.
Thus we must look at an isomorphism of this affine plane curve, by drawing a morphism between $C\cong C^{'} : \{x^2 + y^2 - 1 = 0\}$
Note that any smooth conic is rational, due to being over any algebraically closed field.

\begin{lemma*}{
    \label{1.2.1}
    If $C$ is rational plane curve, and $C\cong C^{'}$, then $C^{'}$ is rational.
}
\end{lemma*}
\begin{lemma*}{
    Any affine smooth cubic that is also infinitely smooth is never rational. And any smooth projective cubic curve is never rational.
}
\end{lemma*}

A counter example of this is if we have genus $\mathfrak{g} = \frac{(d-1)(d-2)}{2}$.
Any $d = 4$ with 3 nodes contains all rational curves without a line to connect all other points.

